\section{Effective escape from fixed affine slices}
\label{as-section}

Let $a,b$ be positive coprime integers, let $V,Q,g$ be integers, and suppose
\[
 F(a,b)=VQ^g,\qquad V\ge1,\quad Q>1,\quad g\ge2.
\]
No power-free or disjoint-support premise on $V$ is needed here.

\begin{theorem}[Three-chart effective exponent bound]\label{as-exponent}
There is an effectively computable constant $C\ge1$, depending only
on the fixed quartic, such that
\[
 \log g\le C(1+\log V+\log(1+M)),\qquad
 M=\min\{a,b,|a-b|\}.
\]
Consequently $1+M\ge e^{-1}g^{1/C}/V$.
\end{theorem}
\begin{proof}
For a nonzero integer $c$, the three seed charts use the two polynomials
\[
 \begin{split}
 P_c(T)&=F(T+c,T)
   =13T^4+26cT^3+20c^2T^2+7c^3T+c^4,\\
 R_c(T)&=F(T,c)=T^4+3cT^3+5c^2T^2+3c^3T+c^4.
 \end{split}
\]
Their discriminants are $117c^{12}$ and their logarithmic coefficient
heights are at most $\log26+4\log(1+|c|)$. Equivalently their roots
are $c/(\alpha_i-1)$ and $c\alpha_i$, respectively, where the
$\alpha_i$ are the four distinct roots of $F(X,1)$. The denominators
do not vanish since $F(1,1)=13$.

Apply B\'erczes--Evertse--Gy\H{o}ry, \emph{Effective results for
Diophantine equations over finitely generated domains}, Theorem~2.3,
in the \href{https://arxiv.org/pdf/1301.7175}{primary paper}.
Use the fixed presentation $\mathbb Z=\mathbb Z[U]/(U)$, with
$r=d=1$ in their notation and polynomial degree four. Their
equation $H(t)=\delta_0y^m$ then has the effective bound
$\log m\le C_0(h+1)$, where $h$ bounds the logarithmic coefficient
and $\delta_0$ heights. Take $\delta_0=V$, $y=Q$ and $m=g$.
The root-of-unity exception is absent because $Q>1$.

The coefficient bounds above give the asserted inequality with
$|c|$ in place of $M$. Apply this with $c=a$, $b$, and $a-b$,
using symmetry for the first choice. All are nonzero: $a=b$ would
force $(a,b)=(1,1)$ and $F=13$, which has no such nontrivial power
divisor. Taking the minimum proves the claim and its rearrangement.
\end{proof}

\begin{theorem}[Effective finiteness on bounded slices]\label{as-finite}
For each fixed $V$ and fixed $H$, only finitely many such primitive
positive seeds have $\min\{a,b,|a-b|\}\le H$. They are effectively
determinable in principle.
\end{theorem}
\begin{proof}
For each nonzero fixed $c$ and fixed $V$, the preceding argument bounds
all possible exponents in $P_c(t)=VQ^g$ or $R_c(t)=VQ^g$ effectively,
even without primitivity. For each of the finitely many exponents,
Theorem~2.2 of the same primary paper supplies an effective finite
list. Its hypotheses hold for these degree-four polynomials with
nonzero discriminant, including exponent two. There are only finitely
many charts and nonzero parameters with $|c|\le H$.
\end{proof}

Thus pure profiles with $g\to\infty$ must escape all three boundaries
at least polynomially in $g$. If $V=g^{o(1)}$ the same holds with
exponent $1/C-o(1)$. This does not exclude general small
$\lambda=\max\{1,\log V\}/g$, which permits much larger residuals.
The consecutive local shadows $(Lk,Lk+1)$ all have fixed difference
$-1$, even when $L$ varies, so only finitely many can be global pure
powers across all $g\ge2$. This is a global restriction beyond their
finite local compatibility; it does not settle moving affine slices.
The two cited effective Diophantine theorems remain external inputs,
not Lean theorems proved in this checkpoint.
